\section{Tito}\label{sec:9.0}

\ruleslabel{General Rule}

\begin{generalrule}
Marshal Tito is represented by a single, backprinted counter that is controlled by the Yugoslav Player at all times. Tito may provide the Yugoslav Player certain advantages in combat and guerrilla reinforcement. However, the Axis Player may attempt to identify, locate, and eliminate Tito.
\end{generalrule}

\ruleslabel{Procedure}

\begin{procedure}
Tito is made available to the Yugoslav Player as part of the Partisan reinforcements of Game-Turn 2 (see 13.92). When initially placed on the map, Tito is always placed with his unidentified counterside showing.
\end{procedure}

\ruleslabel{Cases}

\subsection{Movement of Tito}\label{sec:9.1}

The Tito counter is subject to all normal movement rules and moves as if it
were a Partisan combat unit. However, Tito must end each Yugoslav Movement
Phase stacked with a Partisan unit. Exception: See 9.45.

\subsection{Tito And Combat}\label{sec:9.2}

\case{9.21} If Tito (in any state) is stacked with Partisan units that are attacked, the ratio is shifted one column to the left on the CRT. If Tito is stacked with Partisan units that are attacking, the ratio is shifted one column to the right on the CRT. Tito may provide a maximum of one shift per Combat Phase.

\case{9.22} If Partisan units stacked with Tito are obligated to retreat, Tito must retreat with these units. However, Tito may never be eliminated due to normal combat. If a Partisan stack containing Tito is eliminated, Tito is simply placed on top of any Partisan unit in the same Zone, or if this is impossible, on top of any Partisan unit on the map. If there are no Partisan units currently on the map, Tito is placed aside for the time being, but may return to the game with Partisan replacements (see 13.6).

\case{9.23} Tito may only be eliminated according to the procedure described in Case 9.4. If Tito is eliminated, the Axis Player may shift the ratio one column to the right in all Axis attacks against any Partisan units, and one column to the left in all attacks made by Partisan units. In addition, Chetnik Collaboration die rolls are affected (see 7.22).

\subsection{Tito And Yugoslav Reinforcements And Victory Points}\label{sec:9.3}

\case{9.31} If Tito is present within a Zone during the Guerrilla Reinforcement Phase, the Yugoslav Player receives a variable number of reinforcing Partisan groups (see 7.46).

\case{9.32} During each Game-Turn in which Tito is withdrawn from the map (see 9.43), the number of Partisan groups available by recruitment in Objective Displays or Mountain triangles is halved (round fractions down; see 7.44 and 7.45). If Tito is eliminated, this condition is in effect for the remainder of the game. (Also, see 7.49).

\case{9.33} The Yugoslav Player loses 5 Victory Points for each Game-Turn that Tito is withdrawn from the map (see 9.43). If Tito has been eliminated, the Yugoslav Player immediately loses 25 Victory Points in the ensuing Victory Point Stage and loses 5 Victory Points in each succeeding Victory Point Stage.

\subsection{Identifying And Attacking Tito}\label{sec:9.4}

The Axis Player may make a specific attack against Tito, but only after he has been identified and located. However, no identification, location or attack attempt against Tito may be made after Game-Turn 14.

\case{9.41} The Axis Player may attempt to identify Tito during the Tito Phase of each Game-Turn following the one in which the first Partisan brigade is placed on the map (see 7.6) or in which the Yugoslav Player has first accumulated at least 45 Victory Points (see 14.0). The Axis Player rolls a single die at this time, a 6 indicating that Tito has been identified and a 1 through 5 indicating no effect. If Tito is identified, his counter is flipped over to the appropriate side.

\case{9.42} The Axis Player may attempt to locate Tito during the Tito Phase of each Game-Turn following the one in which he is identified. The Axis Player rolls a single die at this time, a 6 indicating that he has been located and a 1 through 5 indicating no effect. When Tito has been located, the German 501st SS Parachute Battalion becomes available as a reinforcement in the ensuing Axis Reinforcement Phase.

\case{9.43} The Axis Player may attempt to eliminate Tito during any Axis Combat Phase or Segment following the Tito Phase in which Tito is located. However, only one elimination attempt against Tito may be made per game. In order to eliminate Tito, the Axis Player must conduct an attack against a Partisan stack of units containing Tito. The 501st SS Parachute Battalion must participate in this attack. However, before resolving normal combat, the Axis Player rolls a single die. A 6 indicates that Tito is eliminated and a 1 through 5 indicates that Tito is ‘‘withdrawn’’ from the map for a number of full Game-Turns equal to this die roll (see 9.45). After this die roll, the normal combat must be resolved, and the 501st SS must participate (it remains on the map and functions as a normal unit for the remainder of the game).

\case{9.44} As long as the 501st SS participates, an Axis elimination attempt against Tito may be performed during an Anti-Guerrilla Operation.

\case{9.45} If Tito must be withdrawn from the map, his counter is picked up and placed aside for the time being. However, he automatically becomes available as a Partisan reinforcement after a number of full Game-Turns equal to the Axis Player’s elimination attempt die roll (see 9.43). Example: If Tito were attacked on Game-Turn 7 with a die roll of 2, he would become available as a Partisan reinforcement in any Partisan-occupied Objective or Zone Display during the Guerrilla Reinforcement Phase of Game-Turn 10.
